%% ma: L10-B14-0031
%% lop: 10
%% chuong: 5
%% bai: 14
%% dang: 10.14.5
%% loai: DS
%% muc: [NB, NB, TH, VD]
%% dap_an: "ĐĐĐS"
%% nguon: "chuyen-de-cac-so-dac-trung-cua-mau-so-lieu-khong-ghep-nhom-toan-10-chuong-trinh-moi.pdf (GV. Phan Nhật Linh) (P04), Bài 03 Dạng 1 Phần II Câu 8"
%% kien_thuc: "Lớp 10 Bài 13 (số trung bình, tứ phân vị), Bài 14 (khoảng biến thiên, khoảng tứ phân vị, giá trị bất thường)"
%% trang_thai: chinh-thuc
%% ghi_chu: "Tự duyệt theo yêu cầu giáo viên 22/10/2026."
\begin{ex}
Kết quả đo chiều cao (cm) của một nhóm gồm $20$ học sinh được ghi lại như sau:
\begin{center}
\begin{tabular}{|c|c|c|c|c|c|c|c|c|c|}\hline
$150$ & $152$ & $153$ & $154$ & $156$ & $157$ & $159$ & $160$ & $161$ & $162$\\\hline
$164$ & $165$ & $167$ & $168$ & $170$ & $172$ & $174$ & $177$ & $178$ & $181$\\\hline
\end{tabular}
\end{center}
Xét tính đúng sai của các khẳng định sau:
\choiceTF{\True Khoảng biến thiên của mẫu số liệu là $31$}{\True Trung vị của mẫu số liệu là $163$}{\True Khoảng tứ phân vị của mẫu số liệu là $14{,}5$}{Giá trị $181$ là giá trị bất thường của mẫu số liệu}
\loigiai{a) Đúng: $R=181-150=31$.

b) Đúng: $M_e=\dfrac{162+164}2=163$.

c) Đúng: $Q_1=\dfrac{156+157}2=156{,}5$; $Q_3=\dfrac{170+172}2=171$; $\Delta_Q=14{,}5$.

d) Sai: $Q_3+1{,}5\Delta_Q=192{,}75>181$ nên $181$ không phải giá trị bất thường.}
\end{ex}
